\subsection{\cite{1120Arad2012}'s 11-20 money request game}
\label{sec:11-20_app}

The level-$k$ model posits that players differ in how many steps ahead they consider when forming their strategies \citep{LkNagel1995,STAHL1994,STAHL1995}.
The model defines different types of players, from level-$0$ to level-$k$. 
Level-$0$ players use some predefined arbitrary decision rule, while level-$k$ players ($k \geq 1$) best respond assuming others are level-$(k-1)$ reasoners.
Such a model highlights the idea that players' behavior depends not only on their decisions but also on their beliefs about how other people think.

To measure the distribution of level-$k$ thinkers in human populations, AR developed the 11-20 game.
The instructions are:
\begin{quote}
\vspace{-.5cm}
\singlespacing
\emph{You and another player are playing a game in which each player requests an amount of money. The amount must be (an integer) between 11 and 20 shekels.  Each player will receive the amount he requests. A player will receive an additional amount of 20 shekels if he asks for exactly one shekel less than the other player. What amount of money would you request?}
\end{quote}
This \emph{basic} version of the game clearly maps players' levels of reasoning to their choices.
A natural starting point is to choose 20 shekels.
This maximizes the guaranteed payment and provides an obvious starting point for level-0 thinking.
Next, level-1 players, anticipating level-0 players, will choose 20, best respond by requesting 19 shekels to earn the bonus. 
Level-2 players, expecting level-1 behavior, choose 18 shekels, and this pattern continues down to the minimum of 11.
More generally, a player choosing $(20-k)$ shekels plausibly reveals themselves as a level-$k$ thinker.

Interestingly, the 11-20 game does not have a pure strategy Nash equilibrium.
The best response to any choice greater than 11 is to undercut the opponent by 1.
But if the other player chooses 11, also selecting 11 is strictly dominated by every other strategy.
The top row of Table~\ref{tab:AR_results} shows the unique symmetric mixed strategy Nash equilibrium for the game, with most of its density lying between 15 and 17 (levels 3 to 5).

\begin{table}[h!]
\caption{Original results from \cite{1120Arad2012}}
\vspace{-.5cm}
\begin{center}
\begin{minipage}{1 \linewidth}
\label{tab:AR_results}
\footnotesize
\setlength{\tabcolsep}{7pt}
\renewcommand{\arraystretch}{1.25}
\begin{tabular*}{\textwidth}{@{\extracolsep{\fill}}lrrrrrrrrrr}
\hline
\textbf{Shekels Requested} & \textbf{11} & \textbf{12} & \textbf{13} & \textbf{14} & \textbf{15} & \textbf{16} & \textbf{17} & \textbf{18} & \textbf{19} & \textbf{20} \\
\textbf{\boldmath{Level-$k$}} & \boldmath{$L9$} & \boldmath{$L8$} & \boldmath{$L7$} & \boldmath{$L6$} & \boldmath{$L5$} & \boldmath{$L4$} & \boldmath{$L3$} & \boldmath{$L2$} & \boldmath{$L1$} & \boldmath{$L0$} \\
\hline
\emph{Nash Eq. Prediction (\%)} & \emph{0} & \emph{0} & \emph{0} & \emph{0} & \emph{25} & \emph{25} & \emph{20} & \emph{15} & \emph{10} & \emph{5} \\
Basic ($n = 108$) (\%)& 4 & 0 & 3 & 6 & 1 & 6 & 32 & 30 & 12 & 6 \\
Cycle ($n = 72$) (\%)& 1 & 1 & 0 & 1 & 0 & 4 & 10 & 22 & 47 & 13 \\
\hline
\emph{Nash Eq. Prediction (\%)} & \emph{0} & \emph{0} & \emph{0} & \emph{10} & \emph{15} & \emph{15} & \emph{15} & \emph{15} & \emph{15} & \emph{15} \\
Costless ($n = 53$) (\%)& 0 & 4 & 0 & 4 & 4 & 4 & 9 & 21 & 40 & 15 \\
\hline
\end{tabular*}
\end{minipage}
\end{center}
\begin{footnotesize}
\begin{singlespace}
\vspace*{-0.05in}
\emph{Notes:} This table reports the empirical PMFs for three versions of the 11-20 game from \citeauthor{1120Arad2012}. 
In the basic version of this game, two players each request a number between 11-20 shekels and they receive that amount. 
If a player requests exactly one less than their opponent, they win their request plus a 20 shekel bonus. 
In the cycle version, players also receive a 20 shekel bonus if they select 20 and their opponent selects 11. 
The costless version is identical to the basic version, except that players receive 17 shekels if they select any amount other than 20.
The basic and cycle versions of the game share a unique symmetric mixed strategy Nash equilibrium, which is shown in the first row.
The unique mixed strategy Nash equilibrium for the costless version is shown in the 4th row.
\end{singlespace}
\end{footnotesize}
\end{table}

Table~\ref{tab:AR_results} also shows that when AR tested this game on pairs of college students, they deviated significantly from the Nash equilibrium (see row titled Basic).
Most notably, 73\% of participants chose between 17 and 19 shekels, whereas only 45\% of the density for the mixed strategy Nash is on these values.

Finally, Table~\ref{tab:AR_results} provides results from two additional variants of the game from AR.
Both variants of the game still involve two players selecting numbers between 11 and 20.
They differ from the basic version in their payouts (see Appendix~\ref{app:instruct} for the full instructions).
In the \emph{cycle} version, players can earn a bonus of 20 shekels by undercutting the other's request by exactly one or by selecting 20 when the other selects 11. 
This version has the same unique mixed strategy Nash Equilibrium as the basic game because the only affected strategy profiles are (20,11) and (11,20), which are outside the equilibrium support.
However, if there is a distribution of level-$k$ thinkers---as opposed to perfectly rational game-theoretic best-responders---such a change might significantly impact play.

The \emph{costless} version has the same bonus structure as the basic game, but with a different payout for choices below 20.
Here, requesting 20 yields 20 shekels outright, while choosing a lower amount guarantees 17 shekels plus a 20-shekel bonus if the lower request is exactly one less than the other player's. 
It is comparatively ``costless'' to continue undercutting.
This game induces a symmetric mixed strategy Nash equilibrium that is more uniform across the choice set than the basic and cycle versions.

The human responses from these game variants, also from a similar sample of college students, are noticeably shifted towards 18-20 shekels---even with the basic and cycle sharing the same equilibrium.
AR attributes this to the increased salience and payoff of selecting a higher number.
AR concluded that the collective results from these three experiments are best explained by a mix of strategic types consisting of level-0, level-1, level-2, level-3, and random choosers.

The empirical human distributions from these three games comprise our training (the basic version) and validation (the costless and cycle versions) data.
The games are distinct enough such that the human response distributions are different, but still all likely well-explained by similar underlying human choice processes.
